% Document engendré automatiquement par AI Project Reviewer.
% Ne pas modifier à la main : il est réécrit à chaque évaluation.
\documentclass[11pt,a4paper]{article}
\usepackage[utf8]{inputenc}
\usepackage[T1]{fontenc}
\usepackage{array}
\usepackage{longtable}
\usepackage{geometry}
\usepackage{hyperref}
\geometry{margin=2.5cm}

\title{Évaluation de projet}
\date{10/09/2026 à 13h44}
\author{AI Project Reviewer}

\begin{document}
\maketitle

\section*{Contexte de l'analyse}
\begin{tabular}{@{}l>{\raggedright\arraybackslash}p{0.62\textwidth}@{}}
\textbf{Projet} & projet \\
\textbf{Provenance} & dépôt git \\
\textbf{Révision} & 18d93a0 \\
\textbf{Date d'analyse} & 10/09/2026 à 13h44 \\
\textbf{Modèle interrogé} & qwen2.5-coder:1.5b \\
\textbf{Appels au modèle} & 6 \\
\textbf{Durée} & 1 min 02 s \\
\textbf{Configuration} & modèle qwen2.5-coder:1.5b · budget 6000 jetons · 12 fichiers/critère · confiance >= 0.50 · 3 tentative(s) \\
\end{tabular}

\section*{Synthèse}
\begin{longtable}{@{}lrr@{}}
\hline
\textbf{Critère} & \textbf{Note} & \textbf{Maximum} \\
\hline
\endfirsthead
\hline
\textbf{Critère} & \textbf{Note} & \textbf{Maximum} \\
\hline
\endhead
Organisation du projet & 10 & 10 \\
Présence de tests & 4 & 10 \\
Documentation & 10 & 10 \\
Architecture et modularité & 0 & 10 \\
Lisibilité du code & 9 & 10 \\
Respect des principes SOLID & 0 & 10 \\
Pertinence des patrons de conception & 0 & 10 \\
Gestion des erreurs & 10 & 10 \\
Sécurité & 10 & 10 \\
\hline
\textbf{Total} & \textbf{53} & \textbf{90} \\
\hline
\end{longtable}

\begin{center}
\Large\textbf{Note globale : 11,8\,/\,20}
\end{center}

L'évaluation a porté sur 9 critère(s) sur 9, et relevé 0 signalement(s) : 0 de gravité haute, 0 moyenne, 0 basse.

\section{Organisation du projet — 10/10}
Organisation évaluée sur 147 fichier(s).

\subsection*{Points forts}
\begin{itemize}
  \item Le projet déclare sa construction (pom.xml).
  \item Les sources suivent la disposition conventionnelle src/main/java.
  \item Une stratégie de déploiement est fournie (Docker).
  \item Aucun fichier ne dépasse 600 lignes.
\end{itemize}

\section{Présence de tests — 4/10}
Le projet contient 24 fichier(s) de test pour 107 classe(s) de production.

\subsection*{Points forts}
\begin{itemize}
  \item 24 fichier(s) de test pour 107 classe(s) de production, soit un rapport de 0.22.
  \item 2969 lignes de test au total.
\end{itemize}

\subsection*{Points faibles}
\begin{itemize}
  \item Le rapport reste sous 0.50 : une partie des classes n'est pas couverte.
\end{itemize}

\subsection*{Recommandations}
\begin{itemize}
  \item Cibler en priorité les classes sans test qui portent des règles métier.
\end{itemize}

\section{Documentation — 10/10}
Le projet est documenté : README présent, 8 fichier(s) de documentation.

\subsection*{Points forts}
\begin{itemize}
  \item Un README est présent à la racine du projet.
  \item 8 fichier(s) de documentation.
  \item 18 paquet(s) sur 18 disposent d'un package-info.
\end{itemize}

\section{Architecture et modularité — 0/10}
\section{Lisibilité du code — 9/10}
Le code est clair et facile à comprendre.

\subsection*{Points forts}
\begin{itemize}
  \item Le code est organisé en modules clairs.
  \item Les fonctions sont bien définies et explicites.
  \item Les commentaires sont utiles et clarifies le contexte.
\end{itemize}

\subsection*{Points faibles}
\begin{itemize}
  \item Il existe des espaces inutiles.
  \item Il y a une incohérence dans les noms de variables.
  \item Il manque des annotations de documentation.
\end{itemize}

\subsection*{Recommandations}
\begin{itemize}
  \item Enlever les espaces inutiles.
  \item Refaire les noms de variables de manière cohérente.
  \item Ajouter des annotations de documentation.
\end{itemize}

\section{Respect des principes SOLID — 0/10}
\section{Pertinence des patrons de conception — 0/10}
\section{Gestion des erreurs — 10/10}
Le projet a une qualité élevée dans la gestion des erreurs.

\subsection*{Points forts}
\begin{itemize}
  \item Le projet utilise une approche robuste et préventive pour identifier les erreurs.
  \item L'utilisation de bibliothèques et frameworks modernes aide à protéger les applications contre les erreurs.
\end{itemize}

\subsection*{Points faibles}
\begin{itemize}
  \item Il est important de noter que le système de gestion des erreurs pourrait être plus performant si les erreurs étaient traitées à l'origine.
  \item Il est essentiel de mettre en place un mécanisme d'intervention pour corriger les erreurs trouvées.
\end{itemize}

\subsection*{Recommandations}
\begin{itemize}
  \item Le projet devrait être testé avec diverses conditions de performance pour s'assurer qu'il fonctionne correctement.
  \item Il est crucial de mettre en place des mécanismes d'intervention pour corriger les erreurs trouvées.
  \item Le projet devrait être testé avec diverses conditions de performance pour s'assurer qu'il fonctionne correctement.
\end{itemize}

\section{Sécurité — 10/10}
Le projet est considéré comme\_SECURE

\end{document}
